\subsection{Synthetic Task Generation}\label{sec:results-synth}

This section evaluates the synthetic task generator, first at a single level
on the four lesion datasets it targets (ISLES22, Shifts-MS, ATLAS v2.0,
GNC\_705), then under cascade inference on all nine held-out datasets. The
lesion datasets are the hardest in the evaluation: their targets are lesions
on MRI with no counterpart among the training classes.


\paragraph{Synthetic-bank canvas.}
The baseline is the Synth configuration ($p_\text{synth} = 0.5$,
synthetic-bank canvas). It reaches Dice between 0.008 and 0.082 on the four
lesion datasets (\autoref{tab:synth-realhost-ood}). With this canvas, changing
the mix of shape families or $p_\text{synth}$ did not change these results
beyond seed-to-seed variation. Since the bank's images are synthesized from CT
label maps, whereas all four lesion datasets are real MRI, we hypothesized
that the canvas rather than the shapes limits transfer.


\paragraph{Real-host canvas.}
To test this, we continued training from the baseline in two runs that differ
only in the canvas: the synthetic bank (stage 1) or real TotalSegmentator
CT and MRI scans (stage 2). Both use only the scatter-field family, with the same
$p_\text{synth}$ and training budget. The real canvas improves Dice on three of
the four datasets: ISLES22 (0.045 to 0.098), Shifts-MS (0.011 to 0.058), and
GNC\_705 (0.070 to 0.105). ATLAS v2.0 is unchanged (0.033 vs.\ 0.032).
Absolute Dice stays below 0.11 on all four.

Stage 3 extends the real-host canvas to all seven shape families and to all
TotalSegmentator subjects (CT and MRI, all organs), and trains longer. This
raises GNC\_705 to 0.223, while the other three datasets fall slightly below
their stage-2 values. Since stage 3 changes three factors at once, neither the
gain nor the drops can be attributed to a single one.

\begin{table}[t]
  \centering
  \small
  \caption{Single-level Dice on the lesion datasets for successive
  synthetic-task configurations. Stages 1 and 2 differ only in the canvas.
  Stage 3 adds all shape families, all TotalSegmentator subjects as canvas,
  and a longer training budget.}
  \label{tab:synth-realhost-ood}
  \begin{tabular}{@{}lcccc@{}}
    \toprule
    Dataset & Baseline & \makecell{1) Synthetic\\bank} & \makecell{2) Real host,\\scatter} & \makecell{3) Real host,\\all families} \\
    \midrule
    ISLES22    & 0.050 & 0.045 & \textbf{0.098} & 0.087 \\
    Shifts-MS  & 0.008 & 0.011 & \textbf{0.058} & 0.051 \\
    ATLAS v2.0 & 0.028 & \textbf{0.033} & 0.032 & 0.027 \\
    GNC\_705   & 0.082 & 0.070 & 0.105 & \textbf{0.223} \\
    \midrule
    Mean (sample-weighted)$^{\dagger}$ & 0.063 & 0.057 & 0.084 & \textbf{0.153} \\
    \bottomrule
  \end{tabular}
  \par\vspace{0.5em}
  \raggedright\footnotesize
  $^{\dagger}$Weighted by the number of evaluated tasks per dataset (46 for
  Shifts-MS to 1490 for GNC\_705, which carries 61\% of the weight). The
  unweighted mean (0.042, 0.040, 0.073, 0.097) gives the same ranking with a
  smaller stage-3 lead.
\end{table}

\paragraph{Effect on TotalSegmentator.}
\autoref{tab:synth-realhost-indist} reports TotalSegmentator validation Dice
for the same checkpoints. Stages 1 and 2 stay within 0.022 of the baseline in
every row, so the lesion gains of stage 2 do not cost accuracy on the training
distribution. Stage 3 improves all four rows (mean 0.434 to 0.505), which
likely reflects its longer training and larger canvas pool rather than the
synthetic tasks themselves.


\begin{table}[t]
  \centering
  \small
  \caption{TotalSegmentator validation Dice for the checkpoints of
  \autoref{tab:synth-realhost-ood}, by modality and class membership.}
  \label{tab:synth-realhost-indist}
  \begin{tabular}{@{}lcccc@{}}
    \toprule
     & Baseline & \makecell{1) Synthetic\\bank} & \makecell{2) Real host,\\scatter} & \makecell{3) Real host,\\all families} \\
    \midrule
    TotalSeg CT, seen      & 0.479 & 0.476 & 0.474 & \textbf{0.572} \\
    TotalSeg CT, unseen    & 0.326 & 0.348 & 0.337 & \textbf{0.385} \\
    TotalSeg MRI, seen     & 0.505 & 0.500 & 0.498 & \textbf{0.590} \\
    TotalSeg MRI, unseen   & 0.395 & 0.379 & 0.379 & \textbf{0.430} \\
    \midrule
    Mean (sample-weighted)$^{\dagger}$ & 0.434 & 0.435 & 0.430 & \textbf{0.505} \\
    \bottomrule
  \end{tabular}
  \par\vspace{0.5em}
  \raggedright\footnotesize
  $^{\dagger}$Weighted by the number of evaluated tasks per row (155 to 273).
  Unweighted: 0.426, 0.426, 0.422, 0.494. Baseline and stage 1 are tied under
  either weighting.
\end{table}

\paragraph{Cascade inference.}
Synth + Cascade continues the stage-3 checkpoint with 20 epochs of cascade
training.
\autoref{tab:synth-cascade-realhost} compares both with step-matched Medverse.
Synth + Cascade improves over Cascade on MSD Hippocampus (0.303 to 0.518) and
NasalSeg (0.547 to 0.644), is worse on HU\_LWK1 ($-0.028$) and Shifts-MS
($-0.012$), and stays within 0.011 elsewhere. These two gains put it ahead on
the non-lesion mean (0.615, vs.\ 0.552 for Medverse and 0.531 for Cascade),
although it still trails Medverse on both datasets. On the lesion datasets,
the synthetic tasks have little effect under cascade inference: all three
models stay below 0.1, and Medverse is best on Shifts-MS (0.098).

\begin{table}[t]
  \centering
  \small
  \caption{Dice on the held-out datasets under cascade inference. Cascade and
  Synth + Cascade use the same spacings per dataset; Medverse is step-matched
  as in \autoref{tab:cascade-medverse}.}
  \label{tab:synth-cascade-realhost}
  \begin{tabular}{@{}lcccc@{}}
    \toprule
    Dataset & Spacings (mm) & Cascade & Synth + Cascade & Medverse \\
    \midrule
    FLARE22         & $[6, 3, 2.5]$    & \textbf{0.751} & 0.743 & 0.359 \\
    MSD Prostate    & $[3, 1.5, 0.75]$ & \textbf{0.375} & 0.372 & 0.357 \\
    MSD Hippocampus & $[0.5]$          & 0.303 & 0.518 & \textbf{0.686} \\
    HU\_LWK1        & $[6, 3, 1]$      & \textbf{0.152} & 0.124 & 0.020 \\
    NasalSeg        & $[1.5, 0.8]$     & 0.547 & 0.644 & \textbf{0.737} \\
    \midrule
    Mean (sample-weighted)$^{\dagger}$ & & 0.531 & \textbf{0.615} & 0.552 \\
    \midrule
    ISLES22         & $[3, 1.5]$       & 0.069 & \textbf{0.070} & 0.059 \\
    Shifts-MS       & $[3, 1.5]$       & 0.031 & 0.019 & \textbf{0.098} \\
    ATLAS v2.0      & $[3.8, 1.9]$     & 0.026 & \textbf{0.032} & 0.024 \\
    GNC\_705        & $[6, 3, 1.2]$    & 0.024 & \textbf{0.035} & 0.007 \\
    \bottomrule
  \end{tabular}
  \par\vspace{0.5em}
  \raggedright\footnotesize
  $^{\dagger}$Over the five datasets above the second midrule, weighted by the
  number of evaluated tasks (36 for HU\_LWK1 to 650 for FLARE22). Unweighted:
  0.426, 0.480, 0.432 (same ranking). The lesion datasets are excluded, since
  all three models score below 0.1 on all of them.
\end{table}

The stage-3 gain on GNC\_705 (0.223 at a single level) is lost under cascade
inference, where Synth + Cascade reaches 0.035. A likely cause is
localization: small lesions that the coarse level misses fall outside the
region-restricted fine crops and cannot be recovered
(\autoref{sec:methodology-cascade}). \autoref{fig:gnc-localization-fail} shows
one such case: the lesion is a handful of voxels at the edge of the kidney,
and the finest level predicts nothing there (Dice 0.00). Since the single-level results use crops centered on the ground truth, part of this drop reflects the loss of ground-truth localization.

\begin{figure}[t]
  \centering
  \includegraphics[width=\linewidth]{gnc_localization_fail.png}
  \caption{GNC\_705, class hypo\_mask\_r, subject 100014\_v30 (Cascade), with
  the coarse (6\,mm) and middle (3\,mm) levels shown for context. The lesion
  (yellow circle) is a handful of voxels at the edge of the kidney; the finest
  (1.2\,mm) level predicts nothing there (Dice 0.00).}
  \label{fig:gnc-localization-fail}
\end{figure}


\paragraph{Final comparison: Synth + Cascade vs.\ Medverse.}
\autoref{tab:synth-cascade-medverse-final} adds NSD and latency for Synth +
Cascade and step-matched Medverse. NSD agrees with Dice on which model is
better on all nine datasets. On the non-lesion mean, Synth + Cascade leads in
Dice (0.615 vs.\ 0.552) and NSD (0.729 vs.\ 0.631) and is $15\times$ faster
(179 vs.\ 2697\,ms per sample); Medverse is faster only on MSD Hippocampus, the
one single-level dataset. Against Medverse at native depth
(\autoref{tab:medverse-matched-native}), Synth + Cascade wins only FLARE22 and
HU\_LWK1 among the non-lesion datasets, but stays ahead on the sample-weighted
mean (0.615 vs.\ 0.586). \autoref{tab:synth-cascade-medverse-perclass} in the
appendix gives per-class results.

\begin{table}[t]
  \centering
  \small
  \caption{Synth + Cascade (Ours) vs.\ step-matched Medverse (MV): Dice, NSD,
  and mean latency per sample on the nine held-out datasets. Bold marks the
  better value per row for Dice and NSD, and for latency only in the mean
  row.}
  \label{tab:synth-cascade-medverse-final}
  \begin{tabular}{@{}lccc c cc c cc@{}}
    \toprule
    & & \multicolumn{2}{c}{Dice} & & \multicolumn{2}{c}{NSD} & & \multicolumn{2}{c}{Latency (ms)} \\
    \cmidrule(lr){3-4} \cmidrule(lr){6-7} \cmidrule(lr){9-10}
    Dataset & Spacings (mm) & Ours & MV & & Ours & MV & & Ours & MV \\
    \midrule
    FLARE22         & $[6, 3, 2.5]$    & \textbf{0.743} & 0.359 & & \textbf{0.775} & 0.333 & & 274 & 4894 \\
    MSD Prostate    & $[3, 1.5, 0.75]$ & \textbf{0.372} & 0.357 & & \textbf{0.314} & 0.300 & & 278 & 5013 \\
    MSD Hippocampus & $[0.5]$          & 0.518 & \textbf{0.686} & & 0.753 & \textbf{0.905} & & 97  & 50   \\
    HU\_LWK1        & $[6, 3, 1]$      & \textbf{0.124} & 0.020 & & \textbf{0.049} & 0.029 & & 407 & 5336 \\
    NasalSeg        & $[1.5, 0.8]$     & 0.644 & \textbf{0.737} & & 0.791 & \textbf{0.844} & & 106 & 1886 \\
    \midrule
    Mean (sample-weighted)$^{\dagger}$ & & \textbf{0.615} & 0.552 & & \textbf{0.729} & 0.631 & & \textbf{179} & 2697 \\
    \midrule
    ISLES22         & $[3, 1.5]$       & \textbf{0.070} & 0.059 & & \textbf{0.086} & 0.070 & & 145 & 1894 \\
    Shifts-MS       & $[3, 1.5]$       & 0.019 & \textbf{0.098} & & 0.123 & \textbf{0.136} & & 272 & 1911 \\
    ATLAS v2.0      & $[3.8, 1.9]$     & \textbf{0.032} & 0.024 & & \textbf{0.044} & 0.012 & & 151 & 1862 \\
    GNC\_705        & $[6, 3, 1.2]$    & \textbf{0.035} & 0.007 & & \textbf{0.044} & 0.005 & & 225 & 4874 \\
    \bottomrule
  \end{tabular}
  \par\vspace{0.5em}
  \raggedright\footnotesize
  $^{\dagger}$Over the five datasets above the second midrule, weighted by the
  number of evaluated tasks (36 for HU\_LWK1 to 650 for FLARE22). Unweighted:
  Dice 0.480 vs.\ 0.432, NSD 0.536 vs.\ 0.482, latency 232 vs.\ 3436\,ms. Dice
  and NSD are macro-averaged over classes; latency is a mean over samples.
\end{table}